\documentclass[10pt,a4paper]{amsart}
\usepackage[margin=19mm]{geometry}
\usepackage{amsmath,amssymb,mathtools}
\usepackage[scr=rsfs,cal=euler]{mathalfa}
\usepackage{xfrac}
\usepackage[unicode,hidelinks,bookmarks=false]{hyperref}
\newcommand{\ie}{i.e.\ }
\newcommand{\cf}{cf.\ }
\newcommand{\resp}{respectively\ }
\newcommand{\Subsection}{\subsection}
\providecommand{\nobreakdash}{\nobreak\hbox{-}}
\setlength{\emergencystretch}{2em}
\multlinegap0pt
\usepackage{bm}
\usepackage{bbm}
\usepackage{dsfont}
\usepackage{xspace}
\usepackage{cancel}
\usepackage{mathtools}
\usepackage{amsthm}
\makeatletter
\@ifundefined{theorem}{\newtheorem{theorem}{Theorem}}{}
\makeatother
\title[Extension Research of Principal Bundle Constraint System The]{Extension Research of Principal Bundle Constraint System Theory on Ricci-flat K\"ahler Manifolds}
\author{Dongzhe Zheng}
\date{Source: arXiv:2506.11593v1 (CC BY 4.0)}
\begin{document}
\maketitle

\noindent\emph{Source and licence.} Extension Research of Principal Bundle Constraint System Theory on Ricci-flat K\"ahler Manifolds. Version arXiv:2506.11593v1 (CC BY 4.0), distributed under the Creative Commons Attribution 4.0 International licence, as recorded in the arXiv licence metadata for this exact version and shown on the abstract page https://arxiv.org/abs/2506.11593v1. Full rights notice: licenses/arxiv-2506.11593-CC-BY-4.0.txt.\par\smallskip

\noindent\textit{Editorial scope.} Counted unit: the Theorem whose source label is \texttt{None} in the article (source0.tex lines 793--802, source label \texttt{None}), delivered together with the article's own complete original proof of it (source0.tex lines 804--890, one \texttt{proof} environment of 87 lines). No mathematics is written, repaired or re-derived: statement and proof are the article's own text line for line. Required context retained uncounted: none. Every definition, display and label of those lines is retained unchanged. Omitted from this package and not redistributed: 1--792, 803--803, 891--988. Nothing inside a retained excerpt is omitted or altered: every retained body is delivered exactly as the article's lines are written, so no omission receipt of any kind accompanies this package. Citations: the retained excerpts carry no citation command at all, so no bibliography entry of the article is transcribed and no reference is redistributed. Reference adaptation: none was needed and none was made; every reference inside the delivered unit resolves inside it, so no unresolved reference or citation is produced. Printed slips kept literal: no printed word, symbol, hypothesis or number of the counted statement or of its proof was altered. Layout and class: the article's own class is \texttt{extarticle} with packages including inputenc, CJKutf8, geometry, graphicx, booktabs, mathrsfs, bm, xcolor, algorithm, algpseudocode, mathtools, enumitem, tikz-cd, newtxtext; the wrapper is the lane's pinned \texttt{amsart} editorial wrapper at 10pt on A4, which restates the theorem-like environments the retained text needs and adds the article's own macro definitions that the retained text needs (none), with extra wrapper packages (bm, bbm, dsfont, xspace, cancel, mathtools). The delivered numbering is the unit's own. Assets: no retained span contains a figure environment, a graphics-inclusion command, a PDF-page inclusion, a table or a TikZ picture; no image file is copied into this package, the package declares no asset dependency and no image file is redistributed with it. Licence: version arXiv:2506.11593v1 of this article is distributed under the Creative Commons Attribution 4.0 International licence, as recorded in the arXiv licence metadata for this exact version and shown on the abstract page https://arxiv.org/abs/2506.11593v1. A full rights notice is delivered at licenses/arxiv-2506.11593-CC-BY-4.0.txt. Characters: every retained excerpt is pure ASCII, so the delivered engine, XeLaTeX with its default Latin Modern text font, reports no missing character.
\begin{theorem}[Precise Structure Theorem for Spencer Torsion Terms]
In the framework of Spencer spectral sequences, curvature torsion terms $\text{Torsion}^k(\Omega)$ have the following precise structure:

\textbf{Case 1} (Low-dimensional degeneracy case, $\dim M \leq 4$):
$$\text{Torsion}^k(\Omega) = \bigoplus_{\substack{i+2j=k \\ j \geq 1}} H^i_{\text{dR}}(M) \otimes (\text{Sym}^j(\mathfrak{g}^*))^\mathfrak{g} \otimes [\Omega]^j$$

\textbf{Case 2} (General dimension, through spectral sequence filtration structure):
$$\text{Torsion}^k(\Omega) = \frac{\bigcap_{r=2}^N \ker(d_r: E_r^{*,k-*} \to E_r^{*+r, k-*-r+1})}{\sum_{r=2}^N \text{im}(d_r: E_r^{*-r, k-*+r-1} \to E_r^{*,k-*})}$$
where $N$ is the convergence step of the spectral sequence.
\end{theorem}

\begin{proof}
We provide rigorous constructive proofs for both cases.

\textbf{Part I: Proof for Case 1 (Low-dimensional Degeneracy Case)}

When $\dim M \leq 4$, according to Theorem 5.1, the spectral sequence degenerates at the $E_2$ page, i.e., $E_2 = E_\infty$.

\textbf{Substep 1.1: Algebraic Structure Analysis of $E_2$ Page}
$$E_2^{p,q} = H^p_{\text{dR}}(M) \otimes H^q(\mathfrak{g}, \text{Sym}^*(\mathfrak{g}))$$

For semisimple Lie algebra $\mathfrak{g}$, according to the Whitehead lemma and Casimir invariant theory:
$$H^q(\mathfrak{g}, \text{Sym}^k(\mathfrak{g})) = \begin{cases} 
(\text{Sym}^k(\mathfrak{g}^*))^\mathfrak{g} & \text{if } q = 0 \\
0 & \text{if } q \geq 1
\end{cases}$$

where $(\text{Sym}^k(\mathfrak{g}^*))^\mathfrak{g}$ is the space of $\mathfrak{g}$-invariant polynomials, spanned by Casimir invariants of $\mathfrak{g}$.

\textbf{Substep 1.2: Separation of Classical and Torsion Parts}
The total cohomology $H^k_{\text{Spencer}}$ is given by the direct sum of $E_\infty^{p,k-p}$:
$$H^k_{\text{Spencer}} = \bigoplus_{p=0}^k E_\infty^{p,k-p} = \bigoplus_{p=0}^k H^p_{\text{dR}}(M) \otimes (\text{Sym}^{k-p}(\mathfrak{g}^*))^\mathfrak{g}$$

The "classical part" corresponds to contributions from the flat case:
$$H^k_{\text{classical}} = E_\infty^{k,0} = H^k_{\text{dR}}(M) \otimes (\text{Sym}^0(\mathfrak{g}^*))^\mathfrak{g} = H^k_{\text{dR}}(M) \otimes \mathbb{R}$$

The "torsion part" consists of new contributions generated by curvature:
$$\text{Torsion}^k(\Omega) = \bigoplus_{p < k} E_\infty^{p,k-p} = \bigoplus_{p=0}^{k-1} H^p_{\text{dR}}(M) \otimes (\text{Sym}^{k-p}(\mathfrak{g}^*))^\mathfrak{g}$$

\textbf{Substep 1.3: Power Grading Structure of Curvature Classes}
Curvature classes $[\Omega] \in H^2_{\text{dR}}(M, \text{ad} P)$ act on de Rham cohomology through cup products:
$$[\alpha] \mapsto [\alpha] \cup [\Omega]^j$$

This establishes mappings from $H^i_{\text{dR}}(M)$ to $H^{i+2j}_{\text{dR}}(M)$. In the Spencer framework, this corresponds to "jumps":
$$E_2^{i,0} \rightsquigarrow E_2^{i+2j,0}$$

But since the $E_2$ page degenerates, these jumps are actually encoded in the tensor product structure of the $E_2$ page itself.

\textbf{Substep 1.4: Algebraic Derivation of Precise Formula}
Reorganizing the grading structure of torsion terms, grouping by powers of curvature classes. Let $k-p = 2j$ (even case) or $k-p = 2j+1$ (odd case).

For the even case $k-p = 2j$:
$$(\text{Sym}^{2j}(\mathfrak{g}^*))^\mathfrak{g} \sim \text{coefficient space of }[\Omega]^j$$

Therefore:
$$\text{Torsion}^k(\Omega) = \bigoplus_{\substack{i+2j=k \\ j \geq 1}} H^i_{\text{dR}}(M) \otimes (\text{Sym}^j(\mathfrak{g}^*))^\mathfrak{g} \otimes [\Omega]^j$$

where $[\Omega]^j$ denotes the $j$-th power of curvature classes.

\textbf{Part II: Proof for Case 2 (General Dimension Case)}

When $\dim M > 4$, the spectral sequence may not degenerate at the $E_2$ page, requiring consideration of contributions from all higher-order differentials.

\textbf{Substep 2.1: Analysis of Cumulative Effects of Higher-Order Differentials}
Torsion terms consist of cohomology classes that "survive" under all higher-order differentials $d_r$ ($r \geq 2$). Specifically, a class $[\alpha] \in E_r^{p,q}$ belongs to torsion terms if and only if:
\begin{itemize}
\item \textbf{$d_r$-closedness}: $d_r([\alpha]) = 0$ for all $r \geq 2$
\item \textbf{Non-$d_r$-exactness}: $[\alpha] \notin \text{im}(d_r)$ for all $r \geq 2$
\end{itemize}

\textbf{Substep 2.2: Precise Definition of Filtration Structure}
According to spectral sequence theory, elements of the $E_\infty$ page are precisely those classes that "survive to the end":
$$E_\infty^{p,q} = \frac{\bigcap_{r=2}^N \ker(d_r: E_r^{p,q} \to E_r^{p+r,q-r+1})}{\sum_{r=2}^N \text{im}(d_r: E_r^{p-r,q+r-1} \to E_r^{p,q})}$$

For fixed total degree $k = p+q$, torsion terms are defined as:
$$\text{Torsion}^k(\Omega) = \bigoplus_{p=0}^{k-1} E_\infty^{p,k-p}$$

where we exclude the $p = k$ term (classical de Rham part).

\textbf{Substep 2.3: Structural Relationship with Curvature Powers}
Although the precise formula in the general case is quite complex, the "weight" of torsion terms (graded by powers of curvature classes) still satisfies a meaningful decomposition:
$$\text{Torsion}^k(\Omega) = \bigoplus_{j=1}^{\lfloor k/2 \rfloor} \text{Torsion}^k_j(\Omega)$$

where $\text{Torsion}^k_j(\Omega)$ is the "weight $j$" part, generated by $[\Omega]^j$ and its related higher-order differentials.

\textbf{Substep 2.4: Algebraic Foundation of Weight Decomposition}
This structural result is based on the following observations:
\begin{enumerate}
\item The spectral sequence has natural multiplicative structure (cup product)
\item Higher-order differentials $d_r$ act as derivations of this multiplicative structure
\item Powers of curvature classes $[\Omega]$ provide natural grading for this multiplication
\end{enumerate}

Through these algebraic properties, it can be proved that torsion terms of different weights do not mix with each other, yielding the weight decomposition.

\textbf{Conclusion}
We have provided precise structural descriptions of Spencer torsion terms under two different geometric situations. These formulas completely characterize the contribution of non-flat curvature to Spencer cohomology, providing mathematical tools for understanding the topological complexity of constraint systems.
\end{proof}

\end{document}
